\chapter{Coleta e Análise de Resultados}

Este capítulo apresenta os resultados obtidos a partir da aplicação da metodologia proposta, iniciando com a análise detalhada da otimização de parâmetros de pré-processamento, seguida pela avaliação do impacto dos modelos, influência dos \textit{datasets} e comparação das estratégias.

\section{Otimização de Parâmetros de Pré-Processamento em Visão Computacional}

Nesta fase do estudo, realiza-se uma análise detalhada sobre a otimização de parâmetros de pré-processamento em um \textit{pipeline} de aprendizado de máquina aplicado à visão computacional. O objetivo principal consistiu em identificar o valor ótimo do parâmetro de realce de contraste em imagens de um \textit{dataset} específico, visando maximizar o desempenho de um modelo de rede neural convolucional (CNN) em tarefas de classificação. O estudo foi conduzido utilizando uma abordagem sistemática de busca em grade (\textit{grid search}), testando diferentes valores do parâmetro em um conjunto de dados de imagens de isoladores operantes em linhas de transmissão.

\subsection{Metodologia Experimental}

\subsubsection{Conjunto de Dados}
O conjunto de dados utilizado foi o DRNPW (\textit{Dataset for Recognition of Non-ceramic Pollution on Windings}), o qual consiste em imagens de isoladores de linhas de transmissão. Este acervo inclui anotações automáticas para detecção de objetos, totalizando 779 imagens brutas processadas. Após etapas de pré-processamento e redimensionamento para uma resolução fixa de $640 \times 640$ pixels, as imagens foram particionadas em conjuntos de treinamento (468 imagens), validação (87 imagens) e teste (30 imagens). O \textit{dataset} é principalmente projetado para tarefas de detecção e classificação de defeitos estruturais em isoladores, demonstrando grande aptidão para o contexto de manutenção preditiva em infraestrutura elétrica.

\subsubsection{Estratégia de Pré-Processamento}
A técnica de pré-processamento selecionada foi intitulada \textit{ContrastEnhancement}, que recorre à biblioteca PIL (\textit{Python Imaging Library}) para retificar os níveis de contraste das amostras visuais. O método incide sobre o realce de contraste da imagem valendo-se da função \texttt{ImageEnhance.Contrast.enhance(factor)}, na qual a variável \texttt{factor} governa a intensidade da transformação. Valores de \texttt{factor} menores que 1,0 atenuam o contraste, ao passo que frações superiores a 1,0 resultam em intensificação. O hiperparâmetro otimizado foi, de forma exclusiva, o \texttt{factor}, estipulado no domínio dos números em ponto flutuante na faixa abarcada entre 0,5 e 1,611.

\subsubsection{Modelo de Aprendizado de Máquina}
A estratégia de avaliação baseou-se um algoritmo simples de Rede Neural Convolucional (CNN), implementado sob o \textit{framework} TensorFlow/Keras. O projeto estrutural adotado abrange:
\begin{itemize}
    \item Camada de entrada conformada por dimensões $(640, 640, 3)$, condizente com imagens do espectro RGB.
    \item Duas camadas convolucionais dispondo de 32 filtros cada (possuindo \textit{kernels} de $3 \times 3$ e ativação ReLU), sucedidas recursivamente por rotinas de agrupamento máximo (\textit{max-pooling} com janelas de $4 \times 4$).
    \item Camada com operação de achatamento matricial (\textit{Flatten}).
    \item Camada densa composta por 8 neurônios artificiais e ativação ReLU.
    \item Camada de saída dispondo de neurônios correspondentes ao montante de classes do preditor, orientada por uma operação de ativação \textit{softmax}.
\end{itemize}

Durante a compilação optou-se pelo otimizador Adam, atrelado a uma função objetivo de entropia cruzada categórica esparsa (\texttt{sparse\_categorical\_crossentropy}) com monitoramento de desempenho provido pela acurácia. O fluxo de treinamento limitou-se a 5 épocas globais. A avaliação do modelo culminou em inferências sobre o subconjunto de testes extraindo-se as métricas: acurácia, precisão, \textit{recall} e \textit{F1-score}.

\subsubsection{Estratégia de Otimização Paramétrica}
A arquitetura de otimização fundamentou-se no paradigma de Busca em Grade (\textit{Grid Search}), técnica exaustiva que avalia valores uniformemente graduados correspondentes ao hiperparâmetro de interesse. Assinalaram-se 5 instâncias distintas para o \texttt{factor}: 0,5, 0,778, 1,056, 1,333 e 1,611, espaçadas iterativamente valendo-se da função \texttt{numpy.linspace}.

Para cada valor adotado do respectivo \texttt{factor}, o fluxo executou os seguintes passos:
\begin{enumerate}
    \item Pré-processamento sistemático restrito ao valor atual avaliado.
    \item Treinamento estatístico do modelo CNN supervisionado durante as 5 épocas.
    \item Extração autônoma das métricas preditivas oriundas do espécime de teste.
    \item Transcrição final com registro do limiar ideal balizado pela detecção do ponto de máxima acurácia observada.
\end{enumerate}

\subsection{Apoio Teórico do Método Selecionado}

A busca em grade consolida-se através da literatura como um método de eleição para o dimensionamento prático de parâmetros visando o ótimo absoluto perante discretização de instâncias curtas. A operacionalidade atrela-se a princípios fundamentais:
\begin{itemize}
    \item \textbf{Espaço de Busca:} A busca foi realizada de forma uniforme, testando valores entre 0,5 e 1,611 em 5 etapas.
    \item \textbf{Custo Computacional:} Embora a busca em grade possa apresentar ineficiências em cenários com múltiplas variáveis a serem otimizadas, no presente caso, em que há apenas uma variável, tal limitação não se configura como um obstáculo significativo.
    \item \textbf{Atratividade Determinística:} Essa abordagem oferece alta reprodutibilidade científica ao contrário de métodos heurísticos como o \textit{Random Search} ou modelos estocásticos em buscas bayesianas.
\end{itemize}

\subsection{Resultados Reportados e Análise Factual}

O ciclo de otimização contemplou satisfatoriamente as 5 métricas submetidas aos blocos do experimento, refletindo na consolidação dos valores de acurácia observados (vide Tabela \ref{tab:resultados_grid_search}).

\begin{table}[H]
    \centering
    \caption{\label{tab:resultados_grid_search}Resultados da otimização por busca em grade do parâmetro de contraste (\texttt{factor})}
    \vspace{0.2cm}
    \begin{tabular}{ccccc}
        \hline
        \textbf{Fator (\texttt{factor})} & \textbf{Acurácia} & \textbf{Precisão} & \textbf{Recall} & \textbf{F1-Score} \\
        \hline
        0,5     & 0,767    & 0,588    & 0,767  & 0,665 \\
        0,778   & 0,833    & 0,694    & 0,833  & 0,758 \\
        1,056   & 0,800    & 0,640    & 0,800  & 0,711 \\
        1,333   & 0,900    & 0,810    & 0,900  & 0,853 \\
        1,611   & 0,833    & 0,694    & 0,833  & 0,758 \\
        \hline
    \end{tabular}
    \fonte{Autor.}
\end{table}

\indent Mediante minucioso escrutínio dos eventos computados, atestam-se os pertinentes achados analíticos:
\begin{itemize}
    \item \textbf{Otimização Consistente:} O marco otimizado recai sobre o fator de 1,333, onde obtiveram-se índices altamente significativos: 90,0\% de acurácia, 81,0\% de precisão e uma pontuação \textit{F1-score} alocada em 85,3\%. Este evento infere que um acréscimo de contraste moderado maximizou a discernibilidade morfológica do arcabouço convolucional de filtros em prol de detecção das características relevantes ao diagnóstico do problema.
    \item \textbf{Estabilidade nos Vértices Adjacentes:} Diferentemente de cenários com degradação extrema, notou-se que os fatores adjacentes (0,778 e 1,611) mantiveram um desempenho sólido, com acurácia de 83,3\% e \textit{F1-score} em torno de 75,8\%. Isso aponta para uma relativa robustez do modelo em relação a pequenas variações de contraste ao redor do ponto ótimo.
    \item \textbf{Comportamento Côncavo Observável:} A variação do desempenho em relação ao fator de contraste descreve um comportamento característico de pico, onde reduções exageradas no contraste (fator de 0,5) afetam consideravelmente a capacidade preditiva, resultando nas métricas mais baixas observadas no experimento, como acurácia de 76,7\%.
\end{itemize}

Cita-se que o consolidado das medições extraídas foi resguardado visando a reprodutibilidade na trilha \texttt{optimization\_factor\_results.json}.

\subsection{Síntese da Etapa e Encaminhamentos}

As validações empíricas reiteram enfaticamente o papel expressivo da varredura sistemática adotada perfeitamente amoldada ao aperfeiçoamento prévio dos módulos convolucionais. O preditor fixou o \texttt{factor} ótimo estimado como sendo 1,333, com proveitos notáveis na detecção base sobre o DRNPW, assinalados no expressivo valor de aderência quantificado (90,0\% acurácia).

Frente às variações colhidas onde os índices de \textit{score} oscilaram sob constrições ou acréscimos contínuos de contraste visual, sinaliza-se expressamente recomendações iminentes focadas nas rodadas sequenciais englobando:
\begin{itemize}
    \item Verificação de processos autônomos por exploração bayesiana viabilizando eficiências dimensionais mais abrangentes a hiperparâmetros diversificados.
    \item Expansão amostral adensada que viabilize \textit{trials} de refinação granular em torno da região favorável já mapeada no fator de contraste, isolando picos globais exatos.
    \item Sindicância exploratória na arquitetura interna recorrendo a laudos em \textit{logs} transientes ou validações em \textit{checkpointing} a fim de diagnosticar com primor a estocástica de impedimentos paralelos encontrados.
    \item Desdobrar destas simulações elementares as diretrizes vitais à validação robusta para implantações sob vigilância \textit{in loco} em infraestrutura isolante de usinas elétricas convencionais.
\end{itemize}

\section{Impacto dos Modelos no Desempenho dos Processamentos}
Será analisado o impacto dos diferentes modelos de redes neurais no desempenho dos processamentos de imagem.

\section{Influência dos Datasets nos Resultados}
Aqui, será discutida a influência dos diferentes datasets nos resultados dos processamentos de imagem.

\section{Comparação entre Diferentes Estratégias de Processamento}
Serão comparadas as diferentes estratégias de processamento de imagem utilizadas no estudo, destacando as vantagens e desvantagens de cada uma.
